\chapter{Network}
\section{Giao thức Mạng (Network Protocol)}

\section{Giao thức Truyền dẫn Nhị phân (Binary Protocol)}
KBMS sử dụng một giao thức truyền tin nhị phân (Binary Protocol) tùy chỉnh chạy trên nền TCP để đảm bảo tốc độ truyền tải tri thức nhanh nhất và khả năng xử lý đồng thời (concurrency) tốt.

\subsection{Cấu trúc Gói tin (Packet Structure)}

Mỗi thông điệp (Message) gửi đi giữa Client và Server được đóng gói theo định dạng sau:

| Trường | Kích thước | Mô tả |
| :--- | :--- | :--- |
| \textbf{Total Length} | 4 bytes (Int32) | Tổng độ dài các trường tính từ \textit{SessLen} trở đi (bỏ qua MessageType). |
| \textbf{Message Type} | 1 byte | Loại thông điệp (Query, Result, Error, ...). |
| \textbf{SessLen} | 2 bytes (UInt16) | Độ dài của chuỗi Session ID. |
| \textbf{Session ID} | N bytes (UTF-8) | Định danh phiên làm việc của người dùng. |
| \textbf{ReqLen} | 2 bytes (UInt16) | Độ dài của chuỗi Request ID. |
| \textbf{Request ID} | M bytes (UTF-8) | Định danh duy nhất cho mỗi yêu cầu truy vấn. |
| \textbf{Payload} | P bytes (UTF-8) | Nội dung thực tế (ví dụ: câu lệnh KBQL hoặc kết quả JSON). |

---

\subsection{Các loại Thông điệp (Message Types)}

Máy chủ được trang bị \textbf{12 Giao thức Chuẩn} bao quát mọi tình huống Truy vấn, Phân mảnh dữ liệu, Bảo mật và Telemetry Hệ thống (Thiết kế gốc từ `MessageType.cs`):

1.  \textbf{Bảo mật \& Xác thực (Security Auth):}
    *   `LOGIN (0x01)`: Xin cấp quyền phiên làm việc bằng thông tin (User/Pass).
    *   `LOGOUT (0x05)`: Đóng kết nối định danh và trả phiên tài nguyên hệ thống.

2.  \textbf{Truy vấn (Query Routing):}
    *   `QUERY (0x02)`: Chở luồng văn bản gốc (Chuỗi lệnh KBQL).
    *   `RESULT (0x03)`: Gói kết quả suy diễn ở dạng tổng quát (JSON String) .
    *   `ERROR (0x04)`: Payload chứa tọa độ `Line`, `Column` và `Stacktrace` dùng để tô đỏ biên dịch viên của Designer.

3.  \textbf{Kỹ thuật Trả Lô Dữ Liệu Lớn (Tabular Data Streaming - TDS)}
    KBMS xử lý Hàng nghìn Fact mà không làm treo RAM Client bằng luồng C\# phân đoạn kinh điển:
    *   `METADATA (0x06)`: Truyền trước sơ đồ lưới dữ liệu (Tên Field, Bề rộng UI Column).
    *   `ROW (0x07)`: Trình bày Payload theo chuỗi vòng lặp "Lô" (Ví dụ: Batch 100 Rows/Lần gói) gửi liên tục qua luồng Stream TCP. 
    *   `FETCH\_DONE (0x08)`: Cờ khóa hạ màn chặn chuỗi vòng lặp đọc. Tổng hợp độ trễ thực thi CSDL (Execution Time - ms).

4.  \textbf{Hệ Quản trị Trực tiếp (Telemetry System)}
    Dành riêng cho các Client mang quyền hệ thống cấp cao (ROOT Role):
    *   `STATS (0x0A / 10)`: Theo dõi Metrics máy chủ thời gian thực.
    *   `LOGS\_STREAM (0x0B / 11)`: Đăng ký lắng nghe Live Log Socket đang chảy của Hệ thống.
    *   `SESSIONS (0x0C / 12)`: Trích xuất danh bạ quản lý Users Online, kiểm soát giới hạn Concurrent Connections.
    *   `MANAGEMENT\_CMD (0x0D / 13)`: Đỉnh cao phân quyền - Dùng để chạy lệnh thao tác cấp hệ điều hành (Chỉnh sửa Setting, Kill Session, Cấp Role).

---

\subsection{Thuật toán Xử lý Luồng tin (Streaming Algorithm)}

Để tránh hiện tượng treo ứng dụng khi truyền tải lượng tri thức lớn, KBMS áp dụng cơ chế \textbf{Asynchronous Packet Processing}:

1.  \textbf{Header Parsing:} Đọc 4 byte đầu tiên để xác định kích thước gói tin cần nhận.
2.  \textbf{Buffer Allocation:} Cấp phát vùng nhớ vừa đủ cho gói tin đó.
3.  \textbf{Exact Reading:} Sử dụng hàm `ReadExact` để đảm bảo nhận đủ dữ liệu trước khi giải mã (Decode).
4.  \textbf{Payload Dispatching:} Chuyển nội dung UTF-8 vào bộ phân tích cú pháp (Parser) hoặc hiển thị lên UI.

---

\subsection{Bảo mật và Kết nối}

*   \textbf{TCP Keep-Alive:} Duy trì kết nối ổn định giữa KBMS Studio và Server.
*   \textbf{Session Management:} Cho phép người dùng thực hiện nhiều truy vấn đồng thời trên cùng một kết nối mà không bị lẫn lộn dữ liệu (nhờ vào Request ID).
*   \textbf{Error Propagation:} Khi Server gặp lỗi suy diễn, nó sẽ đóng gói toàn bộ Stack Trace và vị trí lỗi vào một `Error Message` đặc biệt để Client hiển thị nháy đỏ trên trình soạn thảo bài toán.

---

\subsection{Kiến trúc Xử lý Chịu tải (Concurrency Threading)}

Để duy trì hiệu năng cao phục vụ cho hàng ngàn điểm kết nối TCP, KBMS Server loại bỏ hoàn toàn mô hình 1-Thread-Per-Client truyền thống, thay vào đó áp dụng cấu trúc xử lý Đa luồng Bất đồng bộ (Async/Await) đỉnh cao của .NET Core:
*   \textbf{Asynchronous I/O Socket:} Cổng nhận (`AcceptTcpClientAsync`) sẽ quăng toàn bộ tiến trình phân tích Packet vào ThreadPool để tránh tắt nghẽn luồng lắng nghe chính (`\_ = HandleClientAsync(client)`).
*   \textbf{Tránh Va Chạm Thread Mạng:} Thuật toán gói dữ liệu sử dụng cờ báo `SemaphoreSlim` để khóa luồng ghi byte, tránh trường hợp Response từ hai yêu cầu bất đồng bộ (Async) va chạm và ghi chồng byte lên nhau trên cùng Network Stream.
*   \textbf{Thread-Safe Sessions:} Việc thu thập và phát tán Request / Log liên quan tới client được lưu trữ tập trung qua lớp `ConnectionManager` bằng bộ từ điển an toàn luồng `ConcurrentDictionary`.

Mỗi khi người dùng gửi một câu lệnh KBQL từ Client, Server sẽ thực hiện một quy trình xử lý đa tầng để hiểu ý định (Intent) và thực thi đúng các phép toán/suy diễn.

\subsection{Vòng đời của một Truy vấn (Query Lifecycle)}

Dưới đây là sơ đồ tuần tự mô tả các bước từ khi Client gửi câu lệnh đến khi nhận được kết quả cuối cùng:

![query\_lifecycle\_v3.png](../assets/diagrams/query\_lifecycle\_v3.png)
\textit{Hình: Sơ đồ Tuần tự Vòng đời Truy vấn (Query Lifecycle)}

---

\subsection{Luồng xử lý Packet (Packet Processing Flow)}

Hệ thống xử lý các byte dữ liệu thô từ Stream và chuyển đổi chúng thành các cấu trúc thông điệp có ý nghĩa:

![packet\_processing\_v3.png](../assets/diagrams/packet\_processing\_v3.png)
*Hình: Sơ đồ\# 09.2. Luồng Xử lý Truy vấn Mạng
*

---

\subsection{Phân tích Từ vựng (Lexer)}

\textbf{Ý tưởng:} Biến chuỗi văn bản (String) thô thành một danh sách các tín hiệu có nghĩa (Tokens).
*   \textbf{Thuật toán Tokenizer:} Quét chuỗi ký tự, nhận diện các từ khóa (`SELECT`, `CREATE`), định danh (tên Concept), toán tử và giá trị số/chuỗi.
*   \textbf{Báo lỗi:} Nếu gặp ký tự lạ, Tokenizer ngay lập tức trả về lỗi kèm vị trí chính xác (Dòng, Cột).

---

\subsection{Phân tích Cú pháp (Parser & AST)}

\textbf{Ý tưởng:} Xây dựng cây cú pháp trừu tượng (Abstract Syntax Tree - AST) từ danh sách Tokens.
*   \textbf{Kỹ thuật:} Sử dụng trình phân tích cú pháp \textbf{Recursive Descent} (Phân tích đi xuống đệ quy) để khớp các Tokens với bộ ngữ pháp của KBQL.
*   \textbf{Cấu trúc AST:} Mỗi câu lệnh được biểu diễn thành một cây các đối tượng. Ví dụ, lệnh `SELECT` sẽ có các nhánh con là `Fields`, `Source`, `WhereCondition`, và `Limit`.

---

\subsection{Thực thi và Cấp phát (Executor & Router)}

Sau khi có AST từ `KBMS.Parser`, gói tin bắt đầu chuyển giao cho Logic Cấu trúc lõi:
1.  \textbf{Chồng Khớp (Binding) \& Routing:} Từng lệnh tách lẻ trong chuỗi văn bản được cắt ra và chuyển qua Router `\_knowledgeManager.Execute(ast, ...)` để thực thi tuần tự.
2.  \textbf{Catalog Check:} Concept được truy vấn có tồn tại không?
3.  \textbf{Schema \& Type Check:} Các biến trong hàm `SELECT` hoặc `CALC()` có hợp lệ về kiểu dữ liệu hay không? Mọi sai phạm đều ngay lập tức bị tóm gọn ở lớp `KBMS.Parser.ParserException`, đính kèm vào đó (Line, Column) và đẩy ngược ra Network Client dưới dạng JSON (`Type: ERROR`) nhằm nhắc người dùng sửa sai.

---

\subsection{Lập kế hoạch Thực thi (Query Execution Plan)}

Hệ thống quyết định cách thức thực hiện truy vấn tối ưu nhất:
*   \textbf{Access Pattern:} Sử dụng Index (B+ Tree) hay quét tuần tự (Full Scan)?
*   \textbf{Join Strategy:} Áp dụng phương pháp JOIN nào để kết hợp các Khái niệm?
*   \textbf{Reasoning Injection:} Nếu truy vấn yêu cầu suy diễn (`INFER` hoặc truy cập các biến suy diễn), hệ thống sẽ tiêm (Inject) thêm các bước gọi vào `ReasoningEngine` trước khi trả kết quả cuối cùng.

---

\subsection{Ví dụ luồng xử lý thực tế}

Câu lệnh: `SELECT name FROM Student WHERE age > 18;`

1.  \textbf{Lexer:} `[SELECT, name, FROM, Student, WHERE, age, >, 18, ;]`
2.  \textbf{Parser:} Tạo đối tượng `SelectStatement` với `Table=Student`, `Filter=(age > 18)`.
3.  \textbf{Binder:} Xác nhận `Student` có trong KB, `name` và `age` là hợp lệ.
4.  \textbf{Executor:}
    *   Gọi `StorageEngine` để lấy dữ liệu từ B+ Tree của `Student`.
    *   Sử dụng `ReasoningEngine` để lọc các bản ghi có `age > 18`.
    *   Chỉ trả về trường `name` cho phía Network Layer.

---

\subsection{Kế hoạch Triển khai Mã Nguồn Backend (Implementation Strategy)}

Gắn liền với kiến trúc C\# của hệ thống, luồng thiết kế cho tầng mạng và Bộ phân tích từ vựng (Network \& Parser) được thực thi theo 4 Giai đoạn định hình thực tế:

\subsubsection{Giai đoạn 1: Quản trị Kết nối (TCP Listener & Async Handle)}
Trực tiếp thiết lập máy chủ `KbmsServer` trên kiến trúc Async/Await I/O. Quăng nhiệm vụ hứng Client đang rỗi vào vùng không chặn (Fire-and-forget qua `\_ = HandleClientAsync()`) nhằm khai phóng ThreadPool. Thiết lập hệ thống quét tệp chết (Cleanup Expired Sessions) nhằm đóng ngắt `Socket` lậu và hoàn trả RAM về máy chủ.

\subsubsection{Giai đoạn 2: Mã hóa Giao thức Nhị phân (Packet Binary Protocol)}
Tiến hành lập trình bộ vỏ `KBMNetwork` để viết định dạng truyền tải tuân thủ tài liệu cấu trúc gói `Protocol.cs`. Viết các thao tác C\# Core cấp thấp `Span<byte>`, `BinaryPrimitives` xử lý `BigEndian` tạc khuôn truyền tin (Đọc đủ Byte Session, RequestID, và Payload).

\subsubsection{Giai đoạn 3: Phân tích Cú Pháp Đệ Quy (Lexer & AST Parser)}
Sử dụng tư duy trình biên dịch tạo module `KBMS.Parser`. Lập trình `Token` và quét nhận dạng Regex cắt chuỗi chữ vô nghĩa thành mảng ký pháp Logic. Xây dựng một cỗ máy phân tích đi xuống (Recursive Descent Parser) quét nối Tokens thành Cây Cú Pháp `AST` (Abstract Syntax Tree), giúp CSDL hiểu ý đồ của con người.

\subsubsection{Giai đoạn 4: Bảo mật & Phiên Giao dịch (Session & Security)}
Quản trị hàng ngàn User kết nối TCP đồng thời thông qua tính năng cấu trúc `ConcurrentDictionary` của lớp `ConnectionManager.cs`. Triển khai tính năng User Auth Protocol (LOGIN Message), xác minh Password mã hóa trên hệ từ điển `UserCatalog`, rà soát quyền sở hữu và hạn chế tấn công. Bắt lỗi AST trả về cấu trúc lỗi (kéo theo Line/Col) trên gói JSON Network gửi tới UI của Designer.

\section{Xác thực Giao thức (Protocol Validation)}

Tầng mạng đảm bảo mọi thông điệp nhị phân được truyền dẫn không sai lệch và có độ trễ thấp nhất.

\subsection{Kiểm thử Gói tin Nhị phân (Binary Packet Testing)}

Sử dụng bộ kịch bản \textbf{`Protocol.cs`} để đóng gói và giải mã mockup:

*   \textbf{Handshake}: Mô phỏng quá trình Client gửi gói tin LOGIN và chờ phản hồi RESULT.
*   \textbf{Streaming}: Kiểm tra khả năng nhận diện dấu hiệu kết thúc luồng dữ liệu (`FETCH\_DONE`).

\subsubsection{Minh chứng Mã nguồn (Handshake):}
\textbf{[Missing Image: code_test_binary_protocol.png]}\\ \textit{Hình 9.1: Kịch bản kiểm thử đóng gói và phân rã gói tin nhị phân.}

\subsubsection{Minh chứng Luồng dữ liệu (Gói tin):}
\textbf{[Missing Image: terminal_test_protocol_hex.png]}\\ \textit{Hình 9.2: Nhật ký gói tin nhị phân (Little-endian) truyền trên Socket.}

\subsection{Kiểm thử Tải Kết nối (Concurrency Proof)}

Mô phỏng 256 người dùng cùng thực thi truy vấn để kiểm tra giới hạn của `Asynchronous Sockets`.

\subsubsection{Minh chứng Kết quả (Concurrency):}
\textbf{[Missing Image: terminal_test_concurrent_clients.png]}\\ \textit{Hình 9.3: Bằng chứng Server xử lý đồng thời nhiều Client mà không nghẽn.}

---

> [!IMPORTANT]
> Toàn bộ giao tiếp giữa Studio/CLI và Server đều dựa trên mô hình Socket bất đồng bộ (Async/Await) của .NET 8, đã được tối ưu hóa qua các bài test tải của `CliServerIntegrationTests`.

